\documentclass[aspectratio=169]{beamer}
\usetheme{Madrid}
\usecolortheme{default}
\usepackage{amsmath}
\usepackage{graphicx}
\usepackage{tikz-cd}
\usepackage{multicol}

\setlength{\parindent}{0pt}
\setlength{\parskip}{1\baselineskip}

\title[Lecture 12]{An Introduction to E-Commerce\\\small with a Focus on Machine Learning, Deep Learning, and Artificial Intelligence}
\subtitle{Lecture 12: MLOps + DataOps for E-Commerce}
\author{Dr. Haitham A. El-Ghareeb}
\institute{Faculty of Computers and Information Sciences \\ Mansoura University \\ Egypt}
\date{\today}

\begin{document}

\begin{frame}
  \titlepage
\end{frame}

\begin{frame}{Session plan (90 minutes)}
  \begin{enumerate}
    \item Why MLOps and DataOps matter in e-commerce
    \item From experimentation to production systems
    \item What MLOps covers
    \item Reproducibility
    \item Pipelines and orchestration
    \item Managerial implications
  \end{enumerate}
\end{frame}

\begin{frame}{Learning objectives}
  By the end of this lecture, you should be able to:
  \begin{itemize}
    \item Explain why production ML in e-commerce requires disciplined operational practices beyond model training.
    \item Distinguish MLOps concerns from DataOps concerns while understanding their overlap in end-to-end commerce systems.
    \item Analyze training-serving skew, feature consistency, monitoring, drift, and deployment risk in production ML workflows.
    \item Design an ML system artifact set that includes reproducibility, observability, and governance requirements.
  \end{itemize}
\end{frame}

\section{Why MLOps and DataOps matter in e-commerce}

\begin{frame}{Why MLOps and DataOps matter in e-commerce}
  \begin{itemize}
    \item How is the model trained repeatedly?
    \item How are features computed consistently?
    \item How is the model deployed safely?
    \item How is degradation detected?
    \item How is the system rolled back when behavior becomes harmful?
  \end{itemize}
\end{frame}

\section{From experimentation to production systems}

\begin{frame}{From experimentation to production systems}
  \begin{itemize}
    \item \textbf{The production gap:} missing or delayed features,; inconsistent identifiers across systems,
    \item \textbf{Why e-commerce is especially operational:} product assortments change,; promotions alter behavior,
  \end{itemize}
\end{frame}

\section{What MLOps covers}

\begin{frame}{What MLOps covers}
  \begin{itemize}
    \item \textbf{Typical MLOps responsibilities:} versioning code, data references, and model artifacts,; training pipelines,
    \item \textbf{What DataOps covers:} data pipeline reliability,; data contracts and validation,
  \end{itemize}
\end{frame}

\section{Reproducibility}

\begin{frame}{Reproducibility}
  \begin{itemize}
    \item \textbf{What should be reproducible:} source code version,; feature definitions,
    \item \textbf{Why this matters in commerce:} new data,; a model change,
  \end{itemize}
\end{frame}

\section{Pipelines and orchestration}

\begin{frame}{Pipelines and orchestration}
  \begin{itemize}
    \item \textbf{Typical pipeline stages:} ingest and validate raw data,; transform features,
    \item \textbf{Batch and streaming coexistence:} batch pipelines for training and aggregate features,; streaming or near-real-time pipelines for serving-time context such as session behavior, current inventory, or recent fraud indicators.
  \end{itemize}
\end{frame}

\section{Feature consistency and training-serving skew}

\begin{frame}{Feature consistency and training-serving skew}
  \begin{itemize}
    \item \textbf{How skew happens:} using different code paths offline and online,; stale or delayed online feature values,
    \item \textbf{Why skew is dangerous:} a fraud model may receive delayed account-age features online,; a recommender may use a stale inventory flag,
    \item \textbf{Feature stores as a concept:} Feature stores are often introduced to reduce this risk by standardizing feature definitions and access patterns across training and serving.
  \end{itemize}
\end{frame}

\section{Model registry and release management}

\begin{frame}{Model registry and release management}
  \begin{itemize}
    \item \textbf{Why a registry matters:} which model versions exist,; which one is approved,
    \item \textbf{Release strategies:} shadow deployment,; canary rollout,
  \end{itemize}
\end{frame}

\section{Monitoring production ML systems}

\begin{frame}{Monitoring production ML systems}
  \begin{itemize}
    \item \textbf{Technical monitoring:} latency,; throughput,
    \item \textbf{Model monitoring:} score distribution shifts,; calibration changes,
    \item \textbf{Business monitoring:} click-through and conversion for recommenders,; fraud-dollar capture and false-positive burden,
  \end{itemize}
\end{frame}

\section{Data drift and concept drift}

\begin{frame}{Data drift and concept drift}
  \begin{itemize}
    \item \textbf{Data drift:} seasonal traffic changes,; new product categories,
    \item \textbf{Concept drift:} the same browsing behavior stops indicating purchase intent after a UI redesign,; fraud patterns evolve so past indicators weaken,
    \item \textbf{Response to drift:} alerting,; retraining,
  \end{itemize}
\end{frame}

\section{CI/CD for ML}

\begin{frame}{CI/CD for ML}
  \begin{itemize}
    \item \textbf{What CI for ML should validate:} unit tests for transformations,; schema and feature checks,
    \item \textbf{What CD for ML should consider:} compatibility with serving infrastructure,; readiness of online features,
  \end{itemize}
\end{frame}

\section{Experiment tracking and model comparison}

\begin{frame}{Experiment tracking and model comparison}
  \begin{itemize}
    \item \textbf{What to track:} dataset reference,; feature set version,
    \item \textbf{Why this matters:} Without structured experiment tracking, teams can neither compare models honestly nor explain why a deployed model was chosen.
  \end{itemize}
\end{frame}

\section{Governance and approvals}

\begin{frame}{Governance and approvals}
  \begin{itemize}
    \item \textbf{Examples of governed models:} fraud and abuse scoring,; dynamic pricing,
    \item \textbf{Typical governance artifacts:} model cards,; risk assessment notes,
  \end{itemize}
\end{frame}

\section{Incident response for ML systems}

\begin{frame}{Incident response for ML systems}
  \begin{itemize}
    \item \textbf{Examples of ML incidents:} a recommender starts surfacing unavailable items,; a fraud model causes a spike in false declines,
    \item \textbf{Preparedness:} clear runbooks,; fallback modes,
  \end{itemize}
\end{frame}

\section{Managerial implications}

\begin{frame}{Managerial implications}
  \begin{itemize}
    \item Organizations often underestimate the operational cost of ML.
    \item The model may be only a small fraction of the total system effort.
    \item Most production pain comes from weak data contracts, poor observability, inconsistent features, and unclear ownership.
  \end{itemize}
\end{frame}

\end{document}
